\section{Limitations}
\label{sec:limitations}

Every limitation below is numbered so that the avenues in Section~\ref{sec:future} pair with it
one to one. Several of them run against our own finding, and those are stated first.

\paragraph{L1. Three components are swept rather than sourced.} The assembled rate is a scenario
in the strict sense. Three of its components have no single sourced value and are drawn
uniformly over stated ranges, so the central figure is a point inside a grid rather than an
estimate with a sampling distribution. We report the share of that grid in which the condition
closes, \result{PassShareOursEasy} percent against the easier labour tax reading on our base, and
state plainly that a share of a uniformly weighted grid is not a probability. The corners of the
space reach further than the pass shares suggest: the analytic maximum on our base is
\result{AnalyticMaxOurs} percent, which exceeds the easier required rate, which is why the claim
this paper makes is about the share of the space rather than about every corner of it.

\paragraph{L2. The base of the taxable shareholder share is load-bearing and contestable.} Ours
is equity in taxable accounts over all US equity outstanding; the Congressional Research Service
uses the same numerator over domestically held stock alone. Ours is the right base for a rate on
a dollar of US surplus whoever holds it, because a foreign holder bears close to no US tax at the
owner level and that dollar is still in the surplus. The choice matters: on the domestic-holder
base the central rate rises to \result{TauKCRSCentral} percent and the share of the grid in which
the condition closes rises to \result{PassShareCRSEasy} percent. The central case falls short on
either base, but how far short depends on the choice, and we report both rather than only ours.

\paragraph{L3. Foreign holders are treated as untaxed at the owner level.} The base construction
treats dividends to foreign holders as bearing no US tax at the owner level, though they bear
withholding tax at treaty rates. Applying a fifteen percent treaty rate raises the assembled rate
to \result{TauKWithholding} percent, and to \result{TauKJointWeakening} percent together with the
filers' own debt share, both still below what the condition requires. The treaty rate is a stated
assumption rather than a measurement, and the effective rate actually borne, which varies by
treaty and with the split between dividends and capital gains realized by foreign holders,
remains open.

\paragraph{L4. The assembled rate is a marginal wedge and the question is about annual revenue.}
Under full expensing the deduction precedes the income, so the rate assembled here is a
present-value wedge over the life of one marginal investment, while replacement asks about tax
collected in a year. The two are comparable only if investment is constant, rates are constant
and the capital stock is stationary. Outside those conditions every comparison in
Section~\ref{sec:fiscal} fails. The transition runs against this paper rather than for it, since
revenue actually collected during an investment boom is below what the assembled rate implies,
but we have not bounded the size of that effect beyond the illustrative path in
Appendix~\ref{app:detail}.

\paragraph{L5. The effective labour tax rate does not vary across the wage distribution.} It is a
single economy-wide rate, and because displacement is concentrated away from the top of the
distribution a flat rate probably overstates the revenue loss from displacing low-paid workers
and understates it from displacing high-paid ones. Using rates by quintile, displacement
concentrated at the bottom would require \result{RequiredBottomQuintile} percent rather than
\result{RequiredTopQuintile} percent at the top. Even that lower requirement exceeds the
assembled rate on either jurisdiction, so the direction is known while the net effect is neither
signed nor bounded.

\paragraph{L6. The pass-through coefficient cannot be sourced.} The boundary in
Section~\ref{sec:fiscal} is drawn over a coefficient for which we have no verified estimate here
or, so far as we located, anywhere. The headline sets it to one and the appendix sweeps it, and
the sweep is the weakest arithmetic in the paper. Its only purpose is directional: incomplete
pass-through can only widen the shortfall, never close it.

\paragraph{L7. The household result is incidence arithmetic, not a forecast.} We move a share of
aggregate wage income to capital income, distribute it by actual ownership and hold debts fixed.
That is incidence under stated conventions. It is not a prediction, not a causal estimate, and
not a general-equilibrium result, and the households it describes would in reality adjust in ways
the arithmetic does not represent.

\paragraph{L8. The classification of the housing agencies moves the holder mapping sevenfold.}
Of the debt owed by households whose capacity falls, the share the federal government holds or
guarantees moves by a factor of about seven between the central classification and the
alternative agency treatment. That is the largest single judgement call affecting the incidence
result, it is a classification choice rather than a measurement, and we have not resolved it.

\paragraph{L9. The bust direction rests on two historical episodes.} The fiscal consequence of a
bust is bracketed by the equity-led case of 2000 to 2002 and the debt-led case of 2007 to 2009,
both verified from published aggregates. Two episodes are not a distribution, the wealth effect
we apply is sourced rather than estimated here, and the output path does not feed bank-channel
damage back into itself.

\paragraph{L10. Off-balance-sheet financing is unmeasured, so the financing verdict is
conditional.} The measured tier is nine registrants' on-balance-sheet facts. Off-balance-sheet
vehicles, hardware-backed lending and securitization contribute nothing to it. The defensible
statement is that on the nine filers' own balance sheets the structure resembles the equity-led
case. The statement that the structure resembles the equity-led case is not defensible, and we do
not make it.
